% Capa
\thispagestyle{empty}
\pagenumbering{roman}

\begin{center}
$institution$
\end{center}


\vspace{5.0cm}

\begin{center}
    \textbf{$title$}
\end{center}

\vspace{5.0cm}

\begin{center}
$author$
\end{center}


\vspace{8cm}


\begin{center}

$city$

$year$
    
\end{center}

\newpage
% Folha de aprovação
\thispagestyle{empty}

\begin{center}
$author-upper$
\end{center}

\vspace{4.0cm}

\begin{center}
    \textbf{$title$}
\end{center}

\vspace{4.0cm}

\begin{adjustwidth}{7.5cm}{0cm}
    \justifying
    $confirm_text$

    \setlength{\parindent}{0pt}

    Orientador: $advisor$.

    Coorientadora: Prof\textsuperscript{a}. Dr\textsuperscript{a}. Fernanda Faria Silva.

\end{adjustwidth}

\vspace{5.0cm}

\begin{center}

$city$

$year$
    
\end{center}

\newpage
% Folha de responsabilidade
\thispagestyle{empty}

\begin{invisible}
.
\end{invisible}

\vspace{28\baselineskip}

\textit{As opiniões expressas neste trabalho são de exclusiva responsabilidade do autor.}


\newpage
\pagestyle{fancy}
% Dedicatoria

\begin{invisible}
.
\end{invisible}

\vspace{28\baselineskip}

\begin{adjustwidth}{7.5cm}{0cm}
    \justifying
    \singlespacing
    $dedication$
\end{adjustwidth}

\newpage
\pagestyle{fancy}

\newcommand{\resumo}{%
\titleformat{\section}[block]{\bfseries\filcenter}{}{1em}{}
}
\resumo

\section*{AGRADECIMENTO}
\justify

\begin{invisible}
.
\end{invisible}

\newpage
% Resumo
\pagestyle{fancy}

\section*{RESUMO}

\begin{spacing}{1.0}
\justify

\setlength{\parindent}{0pt}

Esta dissertação investiga como a digitalização do sistema financeiro brasileiro redefiniu as desigualdades regionais no acesso ao crédito entre 2009 e 2023. Partindo do referencial pós-keynesiano que reconhece o papel ativo da moeda no desenvolvimento regional, a pesquisa articula três eixos: a reconfiguração espacial das redes bancárias, o viés metropolitano da digitalização e os efeitos assimétricos sobre a qualidade do crédito. A estratégia empírica combina a construção de doze indicadores financeiros para 5.570 municípios --- integrando ESTBAN, Pesquisas FEBRABAN e PIB municipal (IBGE) --- com uma análise quasi-experimental em três camadas: (i) datação de quebras estruturais, diferenças-em-diferenças com tratamento contínuo, modelo espacial SLX-DiD, mediação causal e regressão quantílica; (ii) quatro estudos de caso em Minas Gerais (Belo Horizonte, Betim, Juiz de Fora e Almenara); e (iii) cenários contrafactuais e recomendações regulatórias. Os resultados mostram que a digitalização comprimiu o estoque de agências e deteriorou o crédito de longo prazo, com efeitos desigualmente distribuídos no território. O SLX-DiD revela um padrão de deslocamento espacial: a perda de cobertura em um município é compensada por ganhos de centralidade nos vizinhos. A mediação causal indica que a deterioração qualitativa do crédito não decorre do fechamento de agências, mas da migração da carteira remanescente para modalidades curtas e padronizadas. Os estudos de caso mostram que o coeficiente médio dos modelos esconde trajetórias radicalmente distintas, a depender da função financeira que cada município exerce em sua escala de inserção regional. A dissertação conclui que a digitalização financeira não é territorialmente neutra e que políticas de inclusão precisam calibrar instrumentos conforme a heterogeneidade espacial, coordenando regulação bancária, infraestrutura digital e desenvolvimento regional.

\vspace{1em}
\textbf{Palavras-chave:} Centralidade financeira. Digitalização bancária. Desigualdades regionais. Econometria espacial. Inclusão financeira. Brasil.

\setlength{\parindent}{15pt}

\end{spacing}

\newpage
% Abstract
\pagestyle{fancy}

\section*{ABSTRACT}

\begin{spacing}{1.0}
\justify

\setlength{\parindent}{0pt}

This dissertation investigates how the digitalization of the Brazilian financial system has reshaped regional inequalities in credit access between 2009 and 2023. Drawing on the post-Keynesian framework that recognizes money as an active force in regional development, the study pursues three analytical axes: the spatial reconfiguration of bank branch networks, the metropolitan bias of digitalization, and the asymmetric effects on credit quality. The empirical strategy combines the construction of twelve municipal financial indicators for 5,570 municipalities --- integrating ESTBAN, FEBRABAN surveys, and municipal GDP (IBGE) --- with a three-layer quasi-experimental approach: (i) structural break dating, continuous-treatment difference-in-differences, SLX-DiD spatial models, causal mediation analysis, and panel quantile regression; (ii) four territorial case studies in Minas Gerais (Belo Horizonte, Betim, Juiz de Fora, and Almenara); and (iii) counterfactual scenarios and regulatory recommendations. Results show that digitalization compressed the stock of physical branches and eroded long-term credit, with effects unevenly distributed across municipalities. The SLX-DiD reveals a spatial displacement pattern: branch closures in one municipality are offset by centrality gains in neighboring ones. Causal mediation indicates that credit quality deterioration does not operate through branch closures but through a portfolio shift toward short-term, standardized loan modalities in the branches that remain open. The case studies show that average coefficients conceal trajectories that differ radically across territories, depending on the financial function each municipality performs at its scale of regional insertion. The dissertation concludes that financial digitalization is not territorially neutral and that inclusion policies must calibrate instruments according to spatial heterogeneity, coordinating banking regulation, digital infrastructure, and regional development.

\vspace{1em}
\textbf{Keywords:} Financial centrality. Banking digitalization. Regional inequalities. Spatial econometrics. Financial inclusion. Brazil.

\setlength{\parindent}{15pt}

\end{spacing}

\newpage
% Lista de figuras
\renewcommand{\listfigurename}{LISTA DE FIGURAS}
\pagestyle{fancy}
\listoffigures

\newpage
% Lista de tabelas
\renewcommand{\listtablename}{LISTA DE TABELAS}
\pagestyle{fancy}
\listoftables

\newpage
% Lista de siglas
\begin{spacing}{1.5}
\begin{center}
\textbf{LISTA DE SIGLAS}
\end{center}
\vspace{1em}

\begin{list}{}{%
  \setlength{\labelwidth}{2.5cm}%
  \setlength{\labelsep}{0.5cm}%
  \setlength{\leftmargin}{3cm}%
  \setlength{\itemsep}{0pt}%
  \setlength{\parsep}{0pt}%
  \renewcommand{\makelabel}[1]{#1\hfill}%
}
\item[ABNT] Associação Brasileira de Normas Técnicas
\item[ACME] \textit{Average Causal Mediation Effect} (efeito causal médio de mediação)
\item[ACP] Análise de Componentes Principais
\item[ADE] \textit{Average Direct Effect} (efeito direto médio)
\item[ANATEL] Agência Nacional de Telecomunicações
\item[API] \textit{Application Programming Interface}
\item[ATT] \textit{Average Treatment effect on the Treated} (efeito médio do tratamento sobre os tratados)
\item[BIC] \textit{Bayesian Information Criterion} (critério bayesiano de informação)
\item[CBD] \textit{Central Business District} (distrito central de negócios)
\item[CFI] \textit{City Footprint Index} (índice de pegada urbana)
\item[CNPJ] Cadastro Nacional da Pessoa Jurídica
\item[COSIF] Plano Contábil das Instituições do Sistema Financeiro Nacional
\item[COVID-19] \textit{Coronavirus Disease 2019}
\item[DiD] \textit{Difference-in-Differences} (diferenças-em-diferenças)
\item[DT] Depósitos Totais
\item[DV] Depósitos à Vista
\item[EF] Efeitos Fixos
\item[ESTBAN] Estatística Bancária Mensal por Município
\item[FEBRABAN] Federação Brasileira de Bancos
\item[FIRE] \textit{Finance, Insurance and Real Estate} (finanças, seguros e imobiliário)
\item[GMM] \textit{Generalized Method of Moments} (método dos momentos generalizados)
\item[HHI] \textit{Herfindahl-Hirschman Index}
\item[IB] \textit{Internet Banking}
\item[IBGE] Instituto Brasileiro de Geografia e Estatística
\item[IC] Intervalo de Confiança
\item[ICB] Índice de Concentração Bancária
\item[ICO] Índice de Concentração das Operações de Crédito
\item[IHH] Índice Herfindahl-Hirschman
\item[IPCA] Índice Nacional de Preços ao Consumidor Amplo
\item[IRC] Índice Regional de Crédito
\item[IRD] Índice Regional de Depósitos
\item[IRE] Índice Regional de Empréstimos
\item[IRF] Índice Regional de Financiamentos
\item[IRL] Índice Regional de Lucro Bancário
\item[IRR] Índice Regional de Risco
\item[LAC] \textit{Latin America and Caribbean} (América Latina e Caribe)
\item[MB] \textit{Mobile Banking}
\item[MENA] \textit{Middle East and North Africa} (Oriente Médio e Norte da África)
\item[OPC] Operações de Crédito
\item[OSF] \textit{Open Science Framework}
\item[PCA] \textit{Principal Component Analysis} (análise de componentes principais)
\item[PDI] \textit{Pavement Density Index} (índice de densidade pavimentada)
\item[PECM] Modelo de Correção de Erros em Painel
\item[PIB] Produto Interno Bruto
\item[PLB] Preferência pela Liquidez dos Bancos
\item[PLP] Preferência pela Liquidez do Público
\item[Pix] Sistema de Pagamentos Instantâneos Brasileiro
\item[RAIS] Relação Anual de Informações Sociais
\item[RMBH] Região Metropolitana de Belo Horizonte
\item[RMSP] Região Metropolitana de São Paulo
\item[RSS] \textit{Residual Sum of Squares} (soma dos quadrados dos resíduos)
\item[SAR] \textit{Spatial Autoregressive} (modelo autorregressivo espacial)
\item[SDiD] \textit{Synthetic Difference-in-Differences} (diferenças-em-diferenças sintético)
\item[SDM] \textit{Spatial Durbin Model} (modelo espacial de Durbin)
\item[SEM] \textit{Spatial Error Model} (modelo de erro espacial)
\item[SLX] \textit{Spatial Lag of X} (defasagem espacial dos regressores)
\item[SSA] \textit{Sub-Saharan Africa} (África Subsaariana)
\item[UFMG] Universidade Federal de Minas Gerais
\item[UNICAD] Sistema de Informações sobre Entidades de Interesse do Banco Central
\item[VAB] Valor Adicionado Bruto
\item[VECM] \textit{Vector Error Correction Model} (modelo de correção de erros vetorial)
\item[3SLS] \textit{Three-Stage Least Squares} (mínimos quadrados em três estágios)
\end{list}
\end{spacing}

\newpage
% Sumário
\renewcommand{\contentsname}{SUMÁRIO}
\pagestyle{fancy}
\tableofcontents

\titleformat{\section} % redefinir apenas section
  {\normalfont\large\bfseries} % estilo
  {} % prefixo antes do título
  {0em} % espaçamento
  {} % código adicional
\newpage
\pagenumbering{arabic}
